%! Mode:: "TeX:UTF-8"
%!TEX encoding = UTF-8 Unicode
%!TEX TS-program = xelatex
\documentclass{ctexart}
\usepackage{bbm}
\usepackage{tikz}
\usepackage{amsmath,amssymb,amsthm,color,mathrsfs}
\usepackage{hyperref}
\usepackage{cleveref}
\usepackage{enumitem,anysize}%
\usepackage{expl3}
\marginsize{1in}{1in}{1in}{1in}%
\theoremstyle{remark}
\newtheorem*{claim*}{Claim}
\newtheorem{lemma}{Lemma}

\DeclareMathOperator{\ot}{ordertype}
\DeclareMathOperator{\dom}{dom}
\DeclareMathOperator{\ran}{ran}
\DeclareMathOperator{\field}{field}
\DeclareMathOperator{\Ord}{Ord}

\newcommand{\cc}{\mathfrak{c}}

\def\email#1{{\texttt{#1}}}
\newcommand{\isep}[1][0pt]{\addtolength{\itemsep}{#1}}
\newcounter{problem}
\renewcommand{\theproblem}{\Roman{problem}}
\newenvironment{problem}{\refstepcounter{problem}\noindent\color{blue}\begin{tikzpicture}[baseline]%
\node at (-0.02em,0.3em) {$\mathbb{P}$};%
\node[scale=0.7] at (0.2em,-0.0em) {R};%
\node[scale=0.7] at (0.6em,0.4em) {O};%
\node[scale=0.8] at (1.05em,0.25em) {B};%
\node at (1.55em,0.3em) {L};%
\node[scale=0.7] at (1.75em,0.45em) {E};%
\node at (2.35em,0.3em) {M};%
\end{tikzpicture}\theproblem}{}
\crefname{problem}{Problem}{Problem}
%\renewcommand\theprob{{\Roman{problem}}}
\newcommand\mysolution{\begin{tikzpicture}[baseline]%
\node at (-0.04em,0.3em) {$\mathbb{S}$};%
\node[scale=0.7] at (0.35em,0.4em) {O};%
\node at (0.7em,0.3em) {\textit{L}};%
\node[scale=0.7] at (0.95em,0.4em) {U};%
\node[scale=1.1] at (1.19em,0.32em){T};%
\node[scale=0.85] at (1.4em,0.24em){I};%
\node at (1.9em,0.32em){$\mathcal{O}$};%
\node[scale=0.75] at (2.3em,0.21em){\texttt{N}};%
\end{tikzpicture}}
\newenvironment{solution}{\begin{proof}[\mysolution]}{\end{proof}}

%%%%%%%%
\newcommand{\calL}{\mathcal{L}}

\newcommand\<{\langle}
\renewcommand\>{\rangle}
\newcommand\eneg{\mathcal{E}_{\neg}}
\newcommand\eto{\mathcal{E}_{\to}}
\newcommand\N{\mathbb{N}}
\newcommand\subini{\subsetneqq_{init}}
\def\to{\rightarrow}
\newcommand{\calA}{\mathcal{A}}
\newcommand{\xor}{\vee}
\newcommand{\bor}{\bigvee}
\newcommand{\band}{\bigwedge}
\newcommand{\xand}{\wedge}
\newcommand{\minus}{\mathbin{\backslash}}
\newcommand{\calF}{\mathcal{F}}
\newcommand\mi[1]{\mathscr{P}(#1)}
\newcommand{\calB}{\mathcal{B}}
\newcommand{\heq}{\mathop{\hat{=}}}
\newcommand{\hneq}{\mathop{\hat{\neq}}}
\newcommand{\calR}{\mathcal{R}}
\newcommand{\hle}{\mathop{\hat{<}}}
\newcommand{\R}{\mathbb{R}}
\newcommand{\calM}{\mathcal{M}}
\newcommand{\calN}{\mathcal{N}}
\newcommand{\frA}{\mathfrak{A}}
\newcommand{\card}{\mathrm{card}}
\newcommand{\frM}{\mathfrak{M}}
\newcommand{\Th}{\mathrm{Th}}
\newcommand{\ecd}{\mathrm{EC}_\Delta}
\newcommand{\Q}{\mathbb{Q}}
\newcommand{\Z}{\mathbb{Z}}
\newcommand{\oto}{\leftrightarrow}
\newcommand{\hin}{\hat{\in}}
\newcommand{\fun}[2]{{}^{#1}#2}
\newcommand{\A}{\mathbbm{A}}
\newcommand{\fra}{\mathfrak{a}}
\newcommand{\frb}{\mathfrak{b}}
\newcommand{\frp}{\mathfrak{p}}
\crefname{enumi}{}{}







\begin{document}
\title{$\mathbb{ALGEBRAIC\ GEOMETRY}$}
\author{白永乐\\ %% 姓名
SID: 202011150087\\ %% 学号
\email{202011150087@mail.bnu.edu.cn}} %% 电邮
\maketitle
\begin{problem}
    Let $R$ be a Abel ring, $\fra$ is an ideal of $R$, and $\sqrt{\fra}:=\{x\in R:\exists n\in\N,x^n\in\fra\}$. Prove that:
    \begin{enumerate}[ref=\theproblem.\arabic*]
        \item\label{it:11} $\sqrt{\mathfrak{a}}$ is ideal.
        \item\label{it:12} $\sqrt{\sqrt{\mathfrak{a}}}=\sqrt{\mathfrak{a}}$.
        \item\label{it:13} $\sqrt{\fra}$ is the smallest radical ideal contain $\mathfrak{a}$. 
        \item\label{it:14} If $\mathfrak{p}$ is prime ideal, then $\mathfrak{p}$ is radical. 
        \item\label{it:15} $\sqrt{\mathfrak{a}}=\bigcap_{\mathfrak{p}\in\mathcal{P}}\mathfrak{p}$, where $\mathcal{P}$ is the set of all prime ideal contains $\fra$. 
    \end{enumerate}
\end{problem}

\begin{solution}
    \begin{enumerate}
        \item $\forall a,b\in\sqrt{\fra},\exists m,n\in\N,a^m,b^n\in \fra$. Consider $a-b$, we have $(a-b)^{m+n}=\sum_{k=0}^{m+n}\binom{m+n}{k}a^kb^{m+n-k}$. Since $k+m+n-k=m+n$, so $k\geq m$ or $m+n-k\geq n$. So $(a-b)^{m+n}\in\fra$ and thus $a-b\in\sqrt{\fra}$. 
        
        $\forall a\in\sqrt{\fra},b\in R,(ab)^n=a^nb^n$. So $ab\in\sqrt{\fra}$. 
        \item  Obviously $\sqrt{\fra}\subset\sqrt{\sqrt{\fra}}$, so only need to prove $\sqrt{\sqrt{\fra}}\subset \sqrt{\fra}$. Consider $a\in \sqrt{\sqrt{\fra}},\exists n\in\N,a^n\in\sqrt{\fra},\exists m\in\N,(a^n)^m\in\fra$. Thus $a^{mn}\in\fra$, so $a\in\sqrt{\fra}$. So $\sqrt{\fra}=\sqrt{\sqrt{\fra}}$. 
        \item Let $\frb$ is a radical ideal contains $\fra$, then $\forall a\in\sqrt{\fra},\exists n\in\N,a^n\in\fra\subset\frb$. 
        Since $\frb$ is radical, we get $a\in\frb$. So $\sqrt{\fra}\subset \frb$. 
        Noting we have proved $\sqrt{\fra}$ is radical in \Cref{it:12}, so it's the smallest. 
        \item $\forall a\in\sqrt{\frp},\exists n\in\N,a^n\in\frp$. Since $\frp$ is prime, so $a\in\frp$. 
        \item From \Cref{it:13} and \Cref{it:14} we get $\sqrt{\fra}\subset \frp$, so we only need to prove $\bigcap_{\frp\in\mathcal{P}}\frp\subset \sqrt{\fra}$. If not, then $\exists a\notin \sqrt{\fra},\forall \frp\in\mathcal{P},a\in\frp$. Let $\mathcal{I}$ is the set of all ideal contains $\fra$ and not contains any of $a^n,n\in\N$. Since $(\mathcal{I},\subset)$ is partial order, and obviously every chain has upper bound(use union), and $\mathcal{I}\neq \emptyset$($\fra\in\mathfrak{I})$. So there is a maximal element in $\mathcal{I}$(by Zorn's lemma). Assume $\mathfrak{q}\in \mathcal{I}$ is maximal element, we will prove $\mathfrak{q}$ is prime ideal. If not, then $\exists x,y\notin\mathfrak{q},xy\in \mathfrak{q}$. Since $\mathfrak{q}$ is maximal, then $(\mathfrak{q},x),(\mathfrak{q},y)$ contians some $a^n$. Assume $a^n=q_1+xt_1,a^m=q_2+yt_2,q_1,q_2\in \mathfrak{q},t_1,t_2\in R$. Then $a^{m+n}=q_1(q_2+yt_2)+q_2xt_1+xyt_1t_2\in \mathfrak{q}$, contradiction with the definition of $\mathcal{I}$! So $\mathfrak{q}\in\mathcal{P}$. But $a\notin \mathfrak{q}$, contradiction with the assumption that $a\in\frp\forall\frp\in\mathcal{P}$! So $\bigcap_{\frp\in\mathcal{P}}\frp\subset \sqrt{\fra}$. 
    \end{enumerate}
\end{solution}

\begin{problem}
    An algebraically field is not finite field. 
\end{problem}
\begin{solution}
    Assume $F$ is a finite, consider $f(x)=\prod_{a\in F}(x-a)+1\in F[x]$, easily prove $f(x)$ has no root in $F$. 
\end{solution}
\begin{problem}
    Let $A=K[x_1,x_2,\cdots x_n]$, and $m_p=(x_1-a_1,\cdots x_n-a_n),p=(a_1,a_2,\cdots a_n)\in\A_K^n$. Then $m$ is max ideal.
\end{problem}
\begin{lemma}
    \label{lem:ker}
    If $f(x_1,x_2,\cdots x_n)\in K[x_1,x_2,\cdots x_n],f(a_1,a_2,\cdots a_n)=0$, then $f=\sum_{k=1}^n (x_k-a_k)f_k(x_1,x_2,\cdots x_n)$. 
\end{lemma}
\begin{proof}
    Use MI to $n$. When $n=1$ it's obvious. If for some certain $n$ it's right, when goes to $n+1$: Let $g(x_1,x_2,\cdots x_n):=f(x_1,x_2,\cdots x_n,a_{n+1})\in K[x_1,x_2,\cdots x_n]$. Then $g(a_1,a_2,\cdots a_n)=0$, so $g(x_1,x_2,\cdots x_n)=\sum_{k=1}^n (x_k-a_k)g_i(x_1,x_2,\cdots x_n)$. Let $h(x_{n+1}):=f(x_1,x_2,\cdots x_{n+1})-g(x_1,x_2,\cdots x_n)\in K[x_1,x_2,\cdots x_n][x_{n+1}]$, then $h(a_{n+1})=0$. So $h(x_{n+1})=(x_{n+1}-a_{n+1})h_1(x_{n+1})$ for some $h_1(x_{n+1})\in K[x_1,x_2,\cdots x_n][x_{n+1}]$. Then $f(x_1,x_2,\cdots x_{n+1})=\sum_{k=1}^{n+1}(x_{i}-a_i)f_i(x_1,x_2,\cdots x_{n+1})$, where $f_k(x_1,x_2,\cdots x_{n+1})=g_k(x_1,x_2,\cdots x_n),k=1,2,\cdots n$, and $f_{n+1}(x_1,x_2,\cdots x_{n+1})=h_1(x_{n+1})$. 
\end{proof}
\begin{solution}
    Obviously $m_p$ is ideal, so we only need to prove it's max. Consider $\phi:K[x_1,x_2,\cdots x_n]\to K,f(x_1,x_2,\cdots x_n)\mapsto f(a_1,a_2,\cdots a_n)$. Obviously it's a homomorphism, consider $\ker \phi$. Obviously $m_p\subset \ker \phi$, now we prove $\ker \phi\subset m_p$. Assume $f\in\ker\phi$, then $f(a_1,a_2,\cdots a_n)=0$. Use \Cref{lem:ker} we get $f\in\ker\phi$. So $m_p=\ker \phi$. So $R/m_p\cong K$ is a field, thus $m_p$ is max ideal. 
\end{solution}

\begin{problem}
    $A\subset B\subset C$ are Abel rings. If $B$ is f.g. $A-$module and $C$ is f.g. $B-$module, then $C$ is f.g. $A-$module, too.  
\end{problem}
\begin{solution}
    Let $\{b_i:i=1,2,\cdots n\}$ is a basis of $B$ over $A$, and $\{c_i:i=1,2,\cdots m\}$ is a basis of $C$ over $B$. Then for $c\in C,\exists x_i\in B$ such that $c=\sum_{i=1}^m x_i c_i$. And $\exists y_{ij}\in A$ such that $x_i=\sum_{j=1}^n y_{ij} b_j$. So $c=\sum_{i=1}^m\sum_{j=1}^n y_{ij} b_jc_i$. That means $\{b_jc_i:j=1,2,\cdots n,i=1,2,\cdots m\}$ is a basis of $C$ over $A$. 
\end{solution}

\begin{problem}
    If $x$ is integral over $A$ then $A[x]$ is f.g. $A-$module. 
\end{problem}
\begin{solution}
    Assume $x^n+\sum_{k=0}^{n-1}-a_kx^k=0,a_k\in A$. Then we prove $\{x^k:k=0,1,\cdots n-1\}$ is a basis of $A[x]$. Only need to prove $x^m,m\in \N$ can be reperesented. Use MI to $m$. When $m\leq n-1$ it's obvious. Assume for certain $m\geq n,\forall k<m,x^k$ can be reperesented, then for $m$, we have $x^m=x^{m-n}x^n=x^{m-n}\sum_{t=0}^{n-1}a_tx^t=\sum_{t=0}^{n-1}a_tx^{t+m-n}$. Since $t+m-n\leq n-1+m-n=m-1<m$, we get $x^k$ can be reperesented, so $\sum_{t=0}^{n-1}a_tx^{t+m-n}$ can be reperesented. i.e., $x^m$ can be reperesented. So $\{x^k:k=0,1,\cdots n-1\}$ is basis. 
\end{solution}

\begin{problem}
    Let $R$ be an integral domain, finitely generated over a field $k$. If $R$ has transcendence degree $n$ over $k$, then there exist elements $x_1, \ldots, x_n \in R$, algebraically independent over $k$, such that $R$ is integrally dependent on the subring $k\left[x_1, \ldots, x_n\right]$ generated by the $x$ 's.
\end{problem}

\begin{solution}
    
\end{solution}

\end{document}
